\chapter{Conclusion and Future Work}

This report evaluated whether input-conditioned MLP weights can consolidate five independently fine-tuned BERT experts into one model while preserving task performance. The implementation merges non-MLP parameters by Task Arithmetic, upscales MLP modules into weight-ensembled modules controlled by per-layer routers, and trains only the routers by Few-Shot Fine-Tuning. Predictions come from a single forward pass with parameter-level combination and a task-specific head.

Router training is effective: with 256 labelled examples per task, the average score rises from 57.4 for Task Arithmetic to $69.9\pm0.1$ and the gain over majority from 18.4\% to 59.1\%, although a clear gap to the experts (78.2) remains. Small budgets ($k=16,64$) are unstable across seeds. Averages also conceal failures, since the full-MLP model stays below the majority on RTE and the intermediate-only model falls below it on MRPC. Relative to five experts, WEMoE saves 26.9\% of storage with full-MLP upscaling and 52.8\% with intermediate-only upscaling. Because the task identifier is known at inference, the case for consolidation under this design depends on the storage constraint justifying the accuracy loss, and is strongest for the intermediate-only variant.

The contributions are an adaptation of WEMoE to BERT with mask-aware routing, a leakage-controlled protocol, and an analysis that uses a majority floor to expose failures that normalised scores conceal. The findings come from one run with one set of experts and should not be generalised beyond this setting.

Future work follows from the limitations. More seeds, including for the experts and the data split, would give more reliable estimates, and a stricter few-shot protocol would cap development data at $k$. The input-independent router experiment would measure how much of the gain requires input dependence, and measuring inference latency would complete the trade-off evaluation.
